
%%%%%%%%%%%%%%%%%%%%%%%%%%%%%%%%%%%%%%%%%%%%%%%%%%%%%%%%%%%%%
\section{模型的评价}

\subsection{模型的优点}
\begin{itemize}[itemindent=2em]
\item \textbf{优点1：方法系统性强。}本文从影响因素识别到预测模型构建，再到评估优化和策略建议，构建了完整的分析体系。各问题之间层层递进、逻辑清晰，形成了一套可复用的零售数据分析方法论。
\item \textbf{优点2：多维度交叉验证。}在影响因素分析中综合了Pearson、Spearman和随机森林三种方法，在模型评估中采用了交叉验证与网格搜索相结合的策略，确保了结论的可靠性。
\item \textbf{优点3：模型稳健性好。}通过灵敏度分析验证了模型对关键参数不敏感，在参数$\pm20\%$范围内$R^2$波动小于$0.02$，表明模型具有较强的泛化能力和实用价值。
\item \textbf{优点4：可操作性强。}将定量分析结果转化为具体的经营策略建议，通过情景分析和弹性分析为企业提供了清晰的资源投入优先级，具有直接的实践指导意义。
\end{itemize}

\subsection{模型的缺点}
\begin{itemize}[itemindent=2em]
\item \textbf{缺点1：}训练集$R^2$高达$0.999$，存在一定程度的过拟合，测试集$R^2=0.766$与训练集差距明显。后续可引入更严格的正则化策略或早停机制进一步缩小这一差距。
\item \textbf{缺点2：}模型依赖较为丰富的历史数据和特征维度，对于缺乏信息化基础的小型零售门店，部分特征（如线上占比、精确客流数据）可能难以获取，模型的通用性受到一定限制。
\end{itemize}
